\begin{proof}[Proof of \cref{obj:thm:universal-second-order-rate}]
Fix \(K\ge 2\) and a nonzero zero-sum contrast \(c\). We prove the five asserted conclusions in the order displayed.

\begin{enumerate}
\item Let
\[
\sigma_c(t)=\frac{\sum_{a\in\mathcal A_K}c_a t_a}{L_c/2},
\qquad t\in\mathcal T_K .
\]
Since \(c\neq 0\), choose \(a_0\) with \(c_{a_0}\neq 0\). For the response type \(t^{a_0}\) with \(t^{a_0}_{a_0}=1\) and \(t^{a_0}_a=0\) for \(a\neq a_0\), the numerator is \(c_{a_0}\), so the finite set of nonzero values \(|\sigma_c(t)|\) is nonempty and has a positive minimum. This is exactly \(\lambda_c\), hence \(0<\lambda_c\). Also, zero-summability gives
\[
c_{a_0}=-\sum_{a\neq a_0}c_a,
\qquad
|c_{a_0}|\le \sum_{a\neq a_0}|c_a|,
\]
and therefore \(2|c_{a_0}|\le L_c\). Thus one admissible normalized nonzero value is at most one, and the minimum satisfies \(\lambda_c\le 1\). Finally,
\[
A_c=\frac{\sqrt{\lambda_c}}2-\frac{\lambda_c}{16}>0,
\qquad
B_c=\frac{7\lambda_c}{16}>0,
\]
because \(0<\lambda_c\le 1\). Hence
\[
\kappa_c=\frac12\min\{A_c,B_c\}>0 .
\]
% lean: lambdaC_pos_le_one, kappaC_pos

\item The first-order rule \(\widehat\tau_{\mathrm{cHT}}^\star\) from \cref{obj:def:first-order-procedure} is controlled first before projection. For the degenerate case \(n=0\), there is a single empty schedule, its target is
\[
\tau_c(\varnothing)=0,
\]
and the centered contrast-weighted sum is the empty sum, hence equals \(0\). The projection fixes \(0\), so the first-order procedure has risk
\[
R_0(\mathcal D^\star,\widehat\tau_{\mathrm{cHT}}^\star;\varnothing)=0.
\]
Since squared-error risks are nonnegative, this gives \(\rho_0(c)=0\). The reciprocal estimate now begins with positive population size, so assume \(n\ge 1\). For \(z=(t_i)_{i=1}^n\), define, for an assignment vector \(A\in\mathcal A_K^n\),
\[
\xi_i(A_i,z)
=
\sum_{a\in S_c}
 c_a\mathbf 1\{A_i=a\}\frac{t_{i,a}-1/2}{q_a^\star},
\qquad
\widetilde\tau_n(A,z)=\frac1n\sum_{i=1}^n\xi_i(A_i,z).
\]
Under the independent \(q^\star\)-design, inverse-probability cancellation and \(\sum_a c_a=0\) give
\[
\mathbb E\{\xi_i(A_i,z)\mid z\}
=
\sum_{a\in\mathcal A_K} c_a t_{i,a},
\qquad
\mathbb E\{\xi_i(A_i,z)^2\mid z\}=C_0(c),
\]
and hence
\[
\mathbb E\{\widetilde\tau_n(A,z)\mid z\}=\tau_c(z),
\qquad
\operatorname{Var}\{\widetilde\tau_n(A,z)\mid z\}\le \frac{C_0(c)}{n}.
\]
The target lies in the projection interval: for each response type \(t\), the positive and negative coefficient masses are both \(L_c/2\), so
\[
-\frac{L_c}{2}
\le
\sum_{a\in\mathcal A_K}c_a t_a
\le
\frac{L_c}{2},
\qquad
-\frac{L_c}{2}\le \tau_c(z)\le \frac{L_c}{2}.
\]
For every \(y\in[-L_c/2,L_c/2]\), interval projection satisfies
\[
\left|\operatorname{proj}_{[-L_c/2,L_c/2]}(x)-y\right|\le |x-y|,
\]
because the projection fixes \(x\) inside the interval and otherwise moves \(x\) to the nearer endpoint between \(x\) and \(y\). Therefore
\[
\left\{\operatorname{proj}_{[-L_c/2,L_c/2]}(\widetilde\tau_n(A,z))-\tau_c(z)\right\}^2
\le
\left\{\widetilde\tau_n(A,z)-\tau_c(z)\right\}^2.
\]
Taking expectations gives
\[
R_n(\mathcal D^\star,\widehat\tau_{\mathrm{cHT}}^\star;z)\le \frac{C_0(c)}{n}
\qquad\text{for every }z\in\mathcal T_K^n,
\]
and taking the infimum over procedures gives \(\rho_n(c)\le C_0(c)/n\) for \(n\ge1\). Together with the zero-unit branch, this proves the first-order envelope for every \(n\).
% lean: minimaxEnvelopeBound_of_contrastWeightedProcedure

\item Since
\[
j_{\mathrm{tail}}\exp\{-j_{\mathrm{tail}}^{1/3}/8\}\longrightarrow 0
\qquad(j_{\mathrm{tail}}\to\infty),
\]
multiplying by \(4\sqrt{\lambda_c}\) gives
\[
4\sqrt{\lambda_c}\,j_{\mathrm{tail}}\exp\{-j_{\mathrm{tail}}^{1/3}/8\}\longrightarrow 0 .
\]
Because \(A_c>0\), there is \(N_0\) such that, for all integers \(j_{\mathrm{tail}}\ge N_0\),
\[
4\sqrt{\lambda_c}\,j_{\mathrm{tail}}\exp\{-j_{\mathrm{tail}}^{1/3}/8\}\le A_c .
\]
Set \(N_c=\max\{N_0,1\}\). Then the stated tail side condition holds for every integer threshold variable at least \(N_c\).
% lean: shrinkage_tail_prefactor_tendsto_zero

For the risk bound, fix \(n\ge N_c\) and a schedule \(z=(t_i)_{i=1}^n\). Under the independent \(q^\star\)-design, define
\[
V_i=\frac{1}{L_c/2}
\sum_{a\in S_c}
c_a\mathbf 1\{A_i=a\}\frac{t_{i,a}-1/2}{q_a^\star},
\qquad
X_n=\frac1n\sum_{i=1}^n V_i,
\]
\[
\Theta_n(z)=\frac1n\sum_{i=1}^n\sigma_c(t_i),
\qquad
Q_n(z)=\sum_{i=1}^n\sigma_c(t_i)^2 .
\]
Then
\[
\mathbb E(X_n\mid z)=\Theta_n(z),
\qquad
\operatorname{Var}(X_n\mid z)=\frac1n-\frac{Q_n(z)}{n^2},
\]
and the definition of \(\lambda_c\) gives
\[
Q_n(z)\ge \lambda_c\sum_{i=1}^n|\sigma_c(t_i)|
\ge \lambda_c n|\Theta_n(z)|.
\]
For the tail estimate, write
\[
\zeta_i(A)=V_i(A)-\sigma_c(t_i),
\qquad
X_n(A)-\Theta_n(z)=\frac1n\sum_{i=1}^n\zeta_i(A).
\]
The product design makes the variables \(\zeta_i\) independent conditional on \(z\). Also \(V_i(A)\in[-1,1]\) and \(\mathbb E(V_i\mid z)=\sigma_c(t_i)\), so \(\zeta_i\) has conditional mean zero and lies in the interval \([-1-\sigma_c(t_i),1-\sigma_c(t_i)]\), whose length is \(2\). Hoeffding's lemma therefore gives, for every real \(\alpha_{\mathrm{mgf}}\),
\[
\mathbb E\!\left[\exp\{\alpha_{\mathrm{mgf}}\zeta_i\}\mid z\right]
\le
\exp\{\alpha_{\mathrm{mgf}}^2/2\}.
\]
Multiplying the moment-generating-function bounds across the independent units yields
\[
\mathbb E\!\left[\exp\!\left\{\alpha_{\mathrm{mgf}}\sum_{i=1}^n\zeta_i\right\}\mid z\right]
\le
\exp\{n\alpha_{\mathrm{mgf}}^2/2\}.
\]
Thus Chernoff's bound, with \(\alpha_{\mathrm{mgf}}=\omega_{\mathrm{tail}}\) for \(\omega_{\mathrm{tail}}>0\) and with the endpoint \(\omega_{\mathrm{tail}}=0\) immediate, gives
\[
\Pr\!\left\{\sum_{i=1}^n\zeta_i\ge n\omega_{\mathrm{tail}}\mid z\right\}
\le \exp\{-n\omega_{\mathrm{tail}}^2/2\}
\qquad(\omega_{\mathrm{tail}}\ge0).
\]
Applying the same argument to \(-\sum_i\zeta_i\) and taking the union of the two one-sided events gives, for every \(\omega_{\mathrm{tail}}\ge0\),
\[
\Pr\{|X_n-\Theta_n(z)|\ge \omega_{\mathrm{tail}}\mid z\}
\le 2\exp\{-n\omega_{\mathrm{tail}}^2/2\}.
\]
% lean: normalizedScore_mean_variance, normalizedTypeScore_sq_lower, normalizedScore_tail_bound

Put
\[
b_n=n^{-1/3},
\qquad
\varepsilon_n=\frac{\sqrt{\lambda_c}}4 b_n,
\qquad
g_n(x)=\max\{-b_n,\min\{x,b_n\}\}.
\]
The clipped-shrinkage estimator is
\[
\widehat\tau_n^{\mathrm{sh}}
=
\frac{L_c}{2}\{X_n-\varepsilon_n g_n(X_n)\},
\]
followed by projection onto \([-L_c/2,L_c/2]\); that projection fixes the displayed value. Indeed, fix a realized assignment and write \(x=X_n\). The bounds
\[
|x|\le 1,\qquad 0\le b_n\le 1,\qquad 0\le\varepsilon_n\le \frac14\le 1
\]
imply
\[
|x-\varepsilon_n g_n(x)|\le 1.
\]
To see this, use the three regions in the definition of \(g_n\). If \(x<-b_n\), then \(g_n(x)=-b_n\) and
\[
-1\le x\le x+\varepsilon_n b_n\le x+b_n<0\le 1.
\]
If \(b_n<x\), then \(g_n(x)=b_n\) and
\[
-1\le 0<x-b_n\le x-\varepsilon_n b_n\le x\le 1.
\]
In the remaining case, \(-b_n\le x\le b_n\), so \(g_n(x)=x\) and
\[
|x-\varepsilon_n g_n(x)|=(1-\varepsilon_n)|x|\le |x|\le 1.
\]
Thus, for every realized assignment,
\[
-\frac{L_c}{2}
\le
\frac{L_c}{2}\{X_n-\varepsilon_n g_n(X_n)\}
\le
\frac{L_c}{2},
\]
because \(L_c/2\ge0\), and the final interval projection is inactive.
% lean: shrunkenClippedScore_abs_le_one
Let
\[
\mathcal C_n(z)=\operatorname{Cov}\{X_n,g_n(X_n)\mid z\}.
\]
Since \(g_n\) is monotone,
\[
\mathcal C_n(z)\ge 0,
\]
and
\[
\mathbb E\!\left[\{X_n-\varepsilon_n g_n(X_n)-\Theta_n(z)\}^2\mid z\right]
\le
\operatorname{Var}(X_n\mid z)-2\varepsilon_n\mathcal C_n(z)+\varepsilon_n^2b_n^2 .
\]
% lean: finiteDesign_cov_monotone_comp_nonneg, finiteDesign_mse_shrinkage_upper

If \(|\Theta_n(z)|\le b_n/2\), define the far-tail event
\[
\mathscr E_{\mathrm{tail}}(z)
=
\left\{A\in\mathcal A_K^n:
|X_n(A)-\Theta_n(z)|\ge b_n/2\right\}.
\]
The tail bound gives
\[
\Pr\{\mathscr E_{\mathrm{tail}}(z)\mid z\}
\le 2\exp\{-n^{1/3}/8\}.
\]
On \(\mathscr E_{\mathrm{tail}}(z)^c\), the local condition gives
\[
|X_n(A)|
\le |X_n(A)-\Theta_n(z)|+|\Theta_n(z)|
\le b_n,
\]
so \(g_n(X_n(A))=X_n(A)\). Since \(|X_n(A)|\le1\) and \(b_n\le1\), clipping also gives the uniform perturbation bound
\[
|g_n(X_n(A))-X_n(A)|\le 1.
\]
Combining the two displays,
\[
|g_n(X_n(A))-X_n(A)|
\le \mathbf 1\{A\in\mathscr E_{\mathrm{tail}}(z)\}.
\]
Moreover \(|X_n(A)-\Theta_n(z)|\le |X_n(A)|+|\Theta_n(z)|\le2\). Since \(\mathbb E(X_n\mid z)=\Theta_n(z)\),
\[
\mathcal C_n(z)
=
\operatorname{Var}(X_n\mid z)
+\mathbb E\!\left[(X_n-\Theta_n(z))\{g_n(X_n)-X_n\}\mid z\right],
\]
and therefore
\[
\begin{aligned}
\mathcal C_n(z)
&\ge \operatorname{Var}(X_n\mid z)
-\mathbb E\!\left[|X_n-\Theta_n(z)|\,|g_n(X_n)-X_n|\mid z\right]\\
&\ge \operatorname{Var}(X_n\mid z)-2\Pr\{\mathscr E_{\mathrm{tail}}(z)\mid z\}\\
&\ge \operatorname{Var}(X_n\mid z)-4\exp\{-n^{1/3}/8\}.
\end{aligned}
\]
Since \(0\le\varepsilon_n\le1/4\),
\[
1-2\varepsilon_n\ge0,
\qquad
Q_n(z)\ge0.
\]
Substituting \(\operatorname{Var}(X_n\mid z)=n^{-1}-Q_n(z)/n^2\) and the covariance lower bound into the preceding MSE inequality gives
\[
\begin{aligned}
&\mathbb E\!\left[\{X_n-\varepsilon_n g_n(X_n)-\Theta_n(z)\}^2\mid z\right] \\
&\qquad\le
\left(1-2\varepsilon_n\right)
\left\{\frac1n-\frac{Q_n(z)}{n^2}\right\}
+8\varepsilon_n\exp\{-n^{1/3}/8\}
+\varepsilon_n^2b_n^2 \\
&\qquad\le
\frac1n-\frac{2\varepsilon_n}{n}
+8\varepsilon_n\exp\{-n^{1/3}/8\}
+\varepsilon_n^2b_n^2,
\end{aligned}
\]
where the last line drops the nonpositive term
\[-(1-2\varepsilon_n)Q_n(z)/n^2\]. Using
\[
\frac{b_n}{n}=n^{-4/3},
\qquad
\varepsilon_n=\frac{\sqrt{\lambda_c}}4 b_n,
\qquad
\varepsilon_n^2b_n^2
=
\frac{\lambda_c}{16}n^{-4/3},
\]
we obtain
\[
\mathbb E\!\left[\{X_n-\varepsilon_n g_n(X_n)-\Theta_n(z)\}^2\mid z\right]
\le
\frac1n-A_c n^{-4/3}
+2\sqrt{\lambda_c}\,b_n\exp\{-n^{1/3}/8\}.
\]
The choice of \(N_c\) implies
\[
2\sqrt{\lambda_c}\,b_n\exp\{-n^{1/3}/8\}
\le
\frac{A_c}{2}n^{-4/3},
\]
so the local case gives
\[
\mathbb E\!\left[\{X_n-\varepsilon_n g_n(X_n)-\Theta_n(z)\}^2\mid z\right]
\le
\frac1n-\frac{A_c}{2}n^{-4/3}
\le
\frac1n-\kappa_c n^{-4/3}.
\]
% lean: shrinkageProcedure_risk_bound_of_tail

If \(|\Theta_n(z)|>b_n/2\), then the lower bound on \(Q_n(z)\) yields
\[
\frac{Q_n(z)}{n^2}\ge \frac{\lambda_c}{2}n^{-4/3}.
\]
Using \(\mathcal C_n(z)\ge 0\) and the same identity for \(\varepsilon_n^2b_n^2\),
\[
\mathbb E\!\left[\{X_n-\varepsilon_n g_n(X_n)-\Theta_n(z)\}^2\mid z\right]
\le
\frac1n-\frac{7\lambda_c}{16}n^{-4/3}
\le
\frac1n-\kappa_c n^{-4/3}.
\]
The two cases cover every schedule. Multiplying by \((L_c/2)^2=C_0(c)\) gives, uniformly in \(z\),
\[
R_n(\mathcal D^\star,\widehat\tau_n^{\mathrm{sh}};z)
\le
C_0(c)\{n^{-1}-\kappa_c n^{-4/3}\}.
\]
Taking the supremum over schedules proves the asserted clipped-shrinkage risk bound for all \(n\ge N_c\).
% lean: shrinkageProcedure_risk_bound_of_tail

\item Since \(\rho_n(c)\) is the minimax infimum, the shrinkage bound implies, for all \(n\ge N_c\),
\[
\rho_n(c)
\le
C_0(c)\{n^{-1}-\kappa_c n^{-4/3}\}.
\]
Thus
\[
d_n(c)=\frac{C_0(c)}n-\rho_n(c)
\ge
C_0(c)\kappa_c n^{-4/3},
\]
and therefore
\[
n^{4/3}d_n(c)\ge C_0(c)\kappa_c
\qquad(n\ge N_c).
\]
On the other hand, \cref{obj:thm:embedded-two-arm-converse} gives, for all \(n\ge 8\),
\[
d_n(c)\le 43C_0(c)n^{-4/3},
\]
so
\[
n^{4/3}d_n(c)\le 43C_0(c)
\qquad(n\ge 8).
\]
The eventual lower and upper bounds imply
\[
C_0(c)\kappa_c
\le
\liminf_{n\to\infty} n^{4/3}d_n(c)
\le
\limsup_{n\to\infty} n^{4/3}d_n(c)
\le
43C_0(c).
\]
% lean: universal_second_order_rate, embedded_two_arm_converse, secondOrderScale_mul_inverse

\item Let \(a_n\) be a positive regularly varying sequence with index \(\beta\), and suppose
\[
a_n d_n(c)\longrightarrow C
\qquad\text{with }C>0.
\]
Define the auxiliary positive sequence
\[
\upsilon_n=
\begin{cases}
1,& n=0,\\
a_n n^{-4/3},& n\ge 1.
\end{cases}
\]
For \(n\ge1\),
\[
\upsilon_n\, n^{4/3}d_n(c)=a_nd_n(c).
\]
The bounds from the previous step and the convergence to \(C\) imply that, eventually,
\[
\frac{C}{2\cdot 43C_0(c)}
\le
\upsilon_n
\le
\frac{2C}{C_0(c)\kappa_c}.
\]
Regular variation at the scale factor \(2\) gives
\[
\frac{a_{2n}}{a_n}\longrightarrow 2^\beta,
\]
and hence
\[
\frac{\upsilon_{2n}}{\upsilon_n}
=
\frac{a_{2n}}{a_n}\,2^{-4/3}
\longrightarrow
2^{\beta-4/3}.
\]
A positive sequence that is eventually bounded above and away from zero can have a dyadic ratio limit only equal to \(1\): indeed, for the dyadic logarithms
\[
\chi_k=\log \upsilon_{2^k},
\]
the increments \(\chi_{k+1}-\chi_k\) converge to the logarithm of the dyadic ratio, while the eventual boundedness of \(\upsilon_{2^k}\) makes \(\chi_k\) eventually bounded; Cesaro averaging of
\[
\frac{\chi_J-\chi_0}{J}
=
\frac1J\sum_{k=0}^{J-1}(\chi_{k+1}-\chi_k)
\]
therefore forces that logarithm to be \(0\). Applying this to \(\upsilon_n\) gives
\[
2^{\beta-4/3}=1.
\]
Since the map \(r\mapsto 2^r\) is strictly increasing,
\[
\beta=\frac43 .
\]
% lean: dyadic_ratio_limit_eq_one, universal_second_order_rate
\end{enumerate}
\end{proof}
